\documentclass[10pt]{article}
\usepackage[T1]{fontenc}
\usepackage{newtxtext,newtxmath}
\usepackage[paperwidth=8.5in,paperheight=11in,margin=0.55in]{geometry}
\usepackage{booktabs,tabularx,siunitx}
\usepackage[font=small,labelfont=bf]{caption}
\usepackage[active,tightpage]{preview}
\setlength\PreviewBorder{10pt}
\sisetup{detect-all}
\setlength{\tabcolsep}{5pt}
\renewcommand{\arraystretch}{1.12}
\pagestyle{empty}
\begin{document}
\begin{preview}\begin{minipage}{\linewidth}
\centering
\captionof{table}{Experiment B protocol: representative source-building scarcity.}
\label{tab:m5-exp-b-setup}
\begin{tabularx}{\linewidth}{p{0.17\linewidth}p{0.23\linewidth}X}
\toprule
Section & Setting & Value \\
\midrule
Design & Scarcity variable & Distinct source buildings, K \\
 & Budget levels & K=10, 20, 50, 100 \\
 & Row allocation & Approximately 500 rows/building; maximum 50k rows \\
 & Building sampling & Representative nested building ladder \\
 & Anomaly rate & Natural anomaly rate retained from selected buildings \\
 & Meter composition & All meter types in selected buildings \\
\addlinespace[2pt]
Data and models & Features & 137 engineered features \\
 & Models & Tree Ensemble; TabPFN (8 estimators) \\
\addlinespace[2pt]
Evaluation & Holdout & Odd-ID buildings; 724 buildings; 10,137,155 rows \\
 & Metric & PR-AUC by meter; equal-weight four-meter macro \\
\addlinespace[2pt]
Reproducibility & Seed and draws & 42; one draw \\
\bottomrule
\end{tabularx}
\end{minipage}\end{preview}
\end{document}
